% Zone „Das kennst du schon“ – aus bank/quadratische-funktionen/zone.jsonl (zusammenbau v0.1)
\einheitenkopf[zone][Kennst du schon]{Das kennst du schon}
\setcounter{aufgabe}{0}
%% TODO Grundvorstellungs-Aufgabe als erste Hauptnummer der Zone (2.8) fehlt in der Bank

%% TODO Titel: Ich-kann-Satz fehlt in der Bank (Platzhalter: Kettenname) – Nr. 1
%% TODO Anweisung: ein Satz über der ersten Teilaufgabe fehlt in der Bank – Nr. 1
\begin{aufgabe}{Quadrieren auch negativer Zahlen und Dezimalzahlen, Quadrat vor Punkt vor Strich: (−7)² = 49, nicht −49; 0,5² = 0,25}
%% TODO Darstellung für schwach fehlt in der Bank (quadratische-funktionen-zone-f1-v1; Abschnitt „Für schwache Schüler“)
\swz{$6^2$ – wie viel?}{}
%% TODO Darstellung für schwach fehlt in der Bank (quadratische-funktionen-zone-f1-v3; Abschnitt „Für schwache Schüler“)
\swz{$1{,}2^2$ – wie viel?}{}
\end{aufgabe}

%% TODO Titel: Ich-kann-Satz fehlt in der Bank (Platzhalter: Kettenname) – Nr. 2
%% TODO Anweisung: ein Satz über der ersten Teilaufgabe fehlt in der Bank – Nr. 2
\begin{aufgabe}{Koordinaten lesen und eintragen, vier Quadranten, (x | y)-Reihenfolge, Kästchenraster mit 0,5-Einteilung}
\swz{Punkt $A$ – Koordinaten?}{\begin{ksys}[xmin=-4,xmax=4,ymin=-4,ymax=4,ablesen]
\punkt{3}{2}{A}
\end{ksys}}
\swz{Punkt $C$ – Koordinaten?}{\begin{ksys}[xmin=-5,xmax=4,ymin=-4,ymax=4,ablesen]
\punkt{-4}{-3}{C}
\end{ksys}}
\end{aufgabe}

%% TODO Titel: Ich-kann-Satz fehlt in der Bank (Platzhalter: Kettenname) – Nr. 3
%% TODO Anweisung: ein Satz über der ersten Teilaufgabe fehlt in der Bank – Nr. 3
\begin{aufgabe}{Lineare Funktion f(x) = m·x + n: Gerade zeichnen, Funktionswert, Punktprobe, Nullstelle}
%% TODO Darstellung für schwach fehlt in der Bank (quadratische-funktionen-zone-f3-v1; Abschnitt „Für schwache Schüler“)
\swz{$g(x) = 2x + 1$ – $g(4)$?}{}
%% TODO Darstellung für schwach fehlt in der Bank (quadratische-funktionen-zone-f3-v3; Abschnitt „Für schwache Schüler“)
\swz{$g(x) = -3x + 2$ – $g(-2)$?}{}
\end{aufgabe}

%% TODO Titel: Ich-kann-Satz fehlt in der Bank (Platzhalter: Kettenname) – Nr. 4
%% TODO Anweisung: ein Satz über der ersten Teilaufgabe fehlt in der Bank – Nr. 4
\begin{aufgabe}{Binomische Formel ausmultiplizieren und zusammenfassen, auch mit Minus in der Klammer}
%% TODO Darstellung für schwach fehlt in der Bank (quadratische-funktionen-zone-f4-v1; Abschnitt „Für schwache Schüler“)
\swz{$(x + 4)^2$ – ausmultipliziert?}{}
%% TODO Darstellung für schwach fehlt in der Bank (quadratische-funktionen-zone-f4-v3; Abschnitt „Für schwache Schüler“)
\swz{$(x - 6)^2$ – ausmultipliziert?}{}
\end{aufgabe}

%% TODO Titel: Ich-kann-Satz fehlt in der Bank (Platzhalter: Kettenname) – Nr. 5
%% TODO Anweisung: ein Satz über der ersten Teilaufgabe fehlt in der Bank – Nr. 5
\begin{aufgabe}{Lineare Gleichung lösen, Terme mit x auf eine Seite bringen, Schreibform mit Strich}
%% TODO Darstellung für schwach fehlt in der Bank (quadratische-funktionen-zone-f5-v1; Abschnitt „Für schwache Schüler“)
\swz{$\text{$3x + 2 = 14$ – $x$?}$}{}
%% TODO Darstellung für schwach fehlt in der Bank (quadratische-funktionen-zone-f5-v3; Abschnitt „Für schwache Schüler“)
\swz{$\text{$2x + 9 = -3$ – $x$?}$}{}
\end{aufgabe}

%% TODO Titel: Ich-kann-Satz fehlt in der Bank (Platzhalter: Kettenname) – Nr. 6
%% TODO Anweisung: ein Satz über der ersten Teilaufgabe fehlt in der Bank – Nr. 6
\begin{aufgabe}{Schnittpunkt zweier Geraden durch Gleichsetzen}
%% TODO Darstellung für schwach fehlt in der Bank (quadratische-funktionen-zone-f6-v1; Abschnitt „Für schwache Schüler“)
\swz{$\text{$g(x) = x + 1$ und $h(x) = -x + 5$ – Schnittpunkt?}$}{}
%% TODO Darstellung für schwach fehlt in der Bank (quadratische-funktionen-zone-f6-v3; Abschnitt „Für schwache Schüler“)
\swz{$\text{$g(x) = 3x + 4$ und $h(x) = x - 2$ – Schnittpunkt?}$}{}
\end{aufgabe}

%% TODO Titel: Ich-kann-Satz fehlt in der Bank (Platzhalter: Kettenname) – Nr. 7
%% TODO Anweisung: ein Satz über der ersten Teilaufgabe fehlt in der Bank – Nr. 7
\begin{aufgabe}{Quadratwurzel ziehen, auch mit dem Taschenrechner, Näherungswert runden; aus einer negativen Zahl gibt es keine Wurzel}
%% TODO Darstellung für schwach fehlt in der Bank (quadratische-funktionen-zone-f7-v1; Abschnitt „Für schwache Schüler“)
\swz{$\sqrt{64}$ – wie viel?}{}
%% TODO Darstellung für schwach fehlt in der Bank (quadratische-funktionen-zone-f7-v3; Abschnitt „Für schwache Schüler“)
\swz{$\sqrt{7}$ – auf zwei Stellen gerundet?}{}
\end{aufgabe}

%% TODO Titel: Ich-kann-Satz fehlt in der Bank (Platzhalter: Kettenname) – Nr. 8
%% TODO Anweisung: ein Satz über der ersten Teilaufgabe fehlt in der Bank – Nr. 8
\begin{aufgabe}{Quadratische Gleichung durch Wurzelziehen lösen und Normalform mit der p-q-Formel lösen, beide Lösungen angeben}
%% TODO Darstellung für schwach fehlt in der Bank (quadratische-funktionen-zone-f8-v1; Abschnitt „Für schwache Schüler“)
\swz{$\text{$x^2 = 36$ – Lösungen?}$}{}
%% TODO Darstellung für schwach fehlt in der Bank (quadratische-funktionen-zone-f8-v3; Abschnitt „Für schwache Schüler“)
\swz{$\text{$x^2 + 4x - 5 = 0$ – Lösungen?}$}{}
\end{aufgabe}

%% TODO Titel: Ich-kann-Satz fehlt in der Bank (Platzhalter: Kettenname) – Nr. 9
%% TODO Anweisung: ein Satz über der ersten Teilaufgabe fehlt in der Bank – Nr. 9
\begin{aufgabe}{Binomische Formel ausmultiplizieren und zusammenfassen, auch mit Minus in der Klammer – Fehler finden}
\swz[3]{Tim rechnet so: \rechnung{(x - 5)^2 &= x^2 - 25} Finde den Fehler und rechne richtig.}{}
\end{aufgabe}

%% TODO Titel: Ich-kann-Satz fehlt in der Bank (Platzhalter: Kettenname) – Nr. 10
%% TODO Anweisung: ein Satz über der ersten Teilaufgabe fehlt in der Bank – Nr. 10
\begin{aufgabe}{Binomische Formel ausmultiplizieren und zusammenfassen, auch mit Minus in der Klammer – selbst rechnen}
%% TODO Darstellung für schwach fehlt in der Bank (quadratische-funktionen-zone-f4-v6; Abschnitt „Für schwache Schüler“)
\swz{$(x + 7)^2$ – ausmultipliziert?}{}
\end{aufgabe}
